\section{从 K2 到下一代 Agent：哪些问题最难}

\subsection{Token Efficiency 与系统约束}
访谈里一个非常明确的判断是：未来大模型竞争不只是 ``谁的 compute 更大''，而是谁能更高效地消化高质量 token。杨植麟把 Muon 优化器、数据改写、合成数据、测试时扩展能力都放在同一张图里，本质上是在回答一个问题：当高质量训练数据越来越稀缺时，怎样让每个 token 产生更高回报。

\begin{importantbox}{为什么 Token Efficiency 比表面上的 Compute Efficiency 更关键}
\begin{itemize}
    \item 真正高价值的数据本来就少，不能指望无限堆数据解决问题；
    \item Agent 任务训练代价更高，因为还要支付 rollout、tool use、环境执行和评估成本；
    \item 如果 token 利用率不高，即便算力更大，也会在高质量数据瓶颈前停下来。
\end{itemize}
\end{importantbox}

\subsection{K2 之后的三道硬门槛}
如果把访谈里的判断压成更工程化的表述，K2 之后至少还有三道门槛没有真正跨过去：第一是 Agent 泛化，第二是评估体系，第三是系统级闭环。

\begin{table}[H]
\centering
\begin{tabular}{p{2.8cm}p{4.2cm}p{4.5cm}}
\toprule
\textbf{门槛} & \textbf{访谈中的判断} & \textbf{为什么难} \\
\midrule
Agent 泛化 & Benchmark 分高不代表真实场景能做对 & 任务分布长尾，训练分布与实际使用不一致 \\
评估体系 & 复杂端到端任务缺少自动验证方式 & 没有稳定 verifier，就很难高质量做 RL \\
系统闭环 & 一方产品更容易把环境与训练绑在一起 & 第三方生态要靠逆向工程拟合模型训练分布 \\
\bottomrule
\end{tabular}
\caption{K2 之后真正困难的部分，不是 ``会不会做演示''，而是如何稳定扩展到真实任务}
\end{table}

\begin{warningbox}{最容易被低估的不是模型能力，而是评估成本}
访谈里对 Benchmark 的批评非常尖锐：Agent 任务不像数学题或有 test case 的代码任务，很多真实工作没有天然标准答案。没有可靠评估，就无法稳定做强化学习；没有强化学习，就很难把长链路能力内化进模型。这是 Agent 迭代速度慢于外界想象的重要原因。
\end{warningbox}

\subsection{本章小结}
如果说 K2 代表了模型能力的一次向前推进，那么下一阶段真正决定天花板的，将是 token 效率、可靠评估和系统化产品闭环。也正因此，杨植麟在整场访谈里始终把 ``模型''、``数据''、``评估''、``产品'' 放在同一张图里讨论。

\subsection{对组织与产品团队的启示}
这场访谈对研究团队之外也有直接启发。如果未来一方产品会越来越重要，那么组织形态也会跟着变化：模型团队不能只负责训练，产品团队也不能只负责包一层 UI，双方都需要理解工具环境、评估闭环和用户真实任务。

\begin{knowledgebox}{为什么一方产品路线会倒逼组织重构}
\begin{itemize}
    \item 做一方产品意味着环境设计、模型训练和产品交互需要同步规划；
    \item 如果组织仍按 ``模型组''、``产品组'' 完全割裂，很多关键反馈无法回流进训练；
    \item Agent 产品越深入工作流，越需要把产品 telemetry、失败案例和训练数据打通。
\end{itemize}
\end{knowledgebox}

